\chapter{Running and deploying}
\label{chap:run}

Two ways to run the system are supported: one Docker Compose stack on a laptop, or the dashboard on Vercel and everything else on Railway. Neither path puts secrets in the repository. Keys live in the environment, in the Clerk setup file on first run, or in \code{llm\_settings} after a user pastes them.

\section{Local stack}

Docker Compose builds one image per service from the root Dockerfile. Migrations run before the API serves traffic. The worker image includes git and Semgrep.

\begin{table}[H]
  \centering
  \caption{Containers in the local stack.}
  \label{tab:containers}
  \begin{tabular}{@{}lp{9.4cm}@{}}
    \toprule
    Container & Role \\
    \midrule
    \code{postgres} & PostgreSQL 16 with pgvector \\
    \code{redis} & Queues and pub/sub \\
    \code{migrate} & Applies Drizzle migrations, then exits \\
    \code{api} & Core API and the webhook, port 4000 \\
    \code{worker} & Review, index, and eval jobs \\
    \code{web} & Dashboard, port 3000 \\
    \code{mcp} & MCP over HTTP, port 4100 \\
    \code{smee} & Forwards GitHub webhooks to the API through a public channel \\
    \bottomrule
  \end{tabular}
\end{table}

\begin{lstlisting}
docker compose up --build
\end{lstlisting}

Open the dashboard at \code{http://localhost:3000}. The first-run screen asks for the two Clerk keys if they are not already in the environment. Then:

\begin{enumerate}
  \item Create the GitHub App from the dashboard. The compose stack has already created a public webhook URL, and the manifest flow points the App at it.
  \item Choose which repositories to install the App on.
  \item On Repositories, turn reviews on. Index the repo if an embedding key is configured.
  \item Open a pull request. The bot reacts, starts the check run, and posts comments. The review page shows each stage as it finishes.
\end{enumerate}

A second free provider is worth configuring. Gemini's free tier is small. An OpenAI-compatible host can be the fallback so a quota error does not fail the review.

To develop the TypeScript itself, start only Postgres, Redis, and the webhook forwarder in Docker, run migrations, then \code{pnpm dev} for the four apps. Point the forwarder at the API on the host, and set the webhook URL to the smee channel.

\code{pnpm eval} sends the sample diffs to a running API and prints recall, precision, and false alarms. The same action exists as Run on the Quality page.

\section{Vercel and Railway}

Production does not use the smee container. GitHub calls the API's public URL directly.

\begin{itemize}
  \item Railway runs Postgres with pgvector, Redis, the API, the worker, and optionally the MCP server. The image target is selected with the \code{SERVICE} variable (\code{api}, \code{worker}, or \code{mcp}).
  \item Vercel runs \code{apps/web}. It needs the public API URL, the MCP URL if used, and the Clerk keys. The Vercel domain is added to the Clerk allow list.
  \item Create a new GitHub App from the deployed dashboard so its webhook URL is the Railway API, not the local smee channel.
\end{itemize}

The API's \code{WEB\_URL} is the dashboard origin, used for redirects and for CORS. After the domains exist, set that variable and the public API URL, then create the GitHub App.

\section{What the demo does not include}

There is no billing, no GitLab or Bitbucket, and no fine-tuned model. Langfuse and Sentry are not wired up. Provider keys saved from the dashboard are stored as plain text, same as the GitHub App private key; that is an accepted demo limitation, called out in the schema. Do not point the demo at a private repository you cannot share with the model provider: free tiers may train on prompts.
